% GENERATED FILE. Edit canonical lessons, then run npm run generate:books.
% canonical-source-hash: fnv1a64:bff5a3dc31cad1ee
% canonical-lessons: GU-C246-parinam, GU-C246-gun, GU-C246-napas-thavu, GU-C246-abhyas, GU-C246-mahenat-karvi

\chapter{Marks and results}
\label{ch:gu-marks-and-results}
\hlchaptermodality{246}

\begin{chapteropening}
\textbf{By the end of this chapter:} \emph{I can talk about marks, results and working hard for them.}

Complete the last lesson of chapter 246: I can talk about marks, results and working hard for them.
\end{chapteropening}

\section[\texorpdfstring{\gu{પરિણામ} (pariṇām)}{પરિણામ (pariṇām)}]{\gu{પરિણામ} (pariṇām) — result, outcome}

\label{lesson:GU-C246-parinam}



\begin{quote}
Before the new one: say the Gujarati for eraser, then the Gujarati for to pronounce.
\end{quote}

\subsection*{You'll want to know: \gu{પરિણામ}}
\textbf{\gu{પરિણામ}} — \emph{pariṇām} — \textquotedblleft{}result, outcome\textquotedblright{}.

\begin{quote}
\textbf{\gu{આજે} \gu{પરિણામ} \gu{આવે} \gu{છે}.} — \emph{āje pariṇām āve che} — \textquotedblleft{}The results come out today.\textquotedblright{}
\end{quote}

A neuter noun.

\begin{cousinweb}
\textbf{\gu{પરિણામ}} (\emph{pariṇām}) is Sanskrit \emph{pariṇāma}, \textquotedblleft{}change, outcome\textquotedblright{}, taken over whole as a tatsama; EtymDB lists it as borrowed from Sanskrit, and CDIAL has no entry for it (EtymDB).
\end{cousinweb}

\begin{grammarlens}[title={What you've built}]
One more word for everyday life at A2.
\end{grammarlens}

\subsection*{Your turn}
\begin{itemize}
\raggedright
  \item \emph{Say it:} \emph{pariṇām}
  \item \emph{Say it:} \emph{pariṇām}, once more
  \item \emph{Say it:} \emph{āje pariṇām āve che} — The results come out today.
  \item \emph{Recall:} say the Gujarati for eraser, then the Gujarati for to pronounce, then say \emph{pariṇām} again
\end{itemize}

\begin{tcolorbox}[breakable,colback=teal!4,colframe=teal!35!black,title={Before you move on}]
What does \gu{પરિણામ} mean? (\textquotedblleft{}Result, outcome\textquotedblright{}.) Say it once more.
\end{tcolorbox}
\section[\texorpdfstring{\gu{ગુણ} (guṇ)}{ગુણ (guṇ)}]{\gu{ગુણ} (guṇ) — marks; quality}

\label{lesson:GU-C246-gun}



\begin{quote}
Before the new one: say the Gujarati for to pronounce, then the Gujarati for result.
\end{quote}

\subsection*{You'll want to know: \gu{ગુણ}}
\textbf{\gu{ગુણ}} — \emph{guṇ} — \textquotedblleft{}marks; quality\textquotedblright{}.

\begin{quote}
\textbf{\gu{તે} \gu{ગણિતમાં} \gu{સારા} \gu{ગુણ} \gu{લાવે} \gu{છે}.} — \emph{te gaṇitmā̃ sārā guṇ lāve che} — \textquotedblleft{}He gets good marks in maths.\textquotedblright{}
\end{quote}

A masculine noun; it also means a good quality.

\begin{cousinweb}
\textbf{\gu{ગુણ}} (\emph{guṇ}) is Sanskrit \emph{guṇa}, \textquotedblleft{}a strand of a cord\textquotedblright{} and then \textquotedblleft{}quality, good quality\textquotedblright{}, taken over as a tatsama (CDIAL 4190).
\end{cousinweb}

\begin{grammarlens}[title={What you've built}]
One more word for everyday life at A2.
\end{grammarlens}

\subsection*{Your turn}
\begin{itemize}
\raggedright
  \item \emph{Say it:} \emph{guṇ}
  \item \emph{Say it:} \emph{guṇ}, once more
  \item \emph{Say it:} \emph{te gaṇitmā̃ sārā guṇ lāve che} — He gets good marks in maths.
  \item \emph{Recall:} say the Gujarati for to pronounce, then the Gujarati for result, then say \emph{guṇ} again
\end{itemize}

\begin{tcolorbox}[breakable,colback=teal!4,colframe=teal!35!black,title={Before you move on}]
What does \gu{ગુણ} mean? (\textquotedblleft{}Marks; quality\textquotedblright{}.) Say it once more.
\end{tcolorbox}
\section[\texorpdfstring{\gu{નાપાસ} \gu{થવું} (nāpās thavũ)}{નાપાસ થવું (nāpās thavũ)}]{\gu{નાપાસ} \gu{થવું} (nāpās thavũ) — to fail an exam}

\label{lesson:GU-C246-napas-thavu}



\begin{quote}
Before the new one: say the Gujarati for result, then the Gujarati for marks.
\end{quote}

\subsection*{You'll want to know: \gu{નાપાસ} \gu{થવું}}
\textbf{\gu{નાપાસ} \gu{થવું}} — \emph{nāpās thavũ} — \textquotedblleft{}to fail an exam\textquotedblright{}.

\begin{quote}
\textbf{\gu{તે} \gu{ગણિતમાં} \gu{નાપાસ} \gu{થાય} \gu{છે}.} — \emph{te gaṇitmā̃ nāpās thāy che} — \textquotedblleft{}He fails in maths.\textquotedblright{}
\end{quote}

\begin{cousinweb}
\textbf{\gu{નાપાસ} \gu{થવું}} (\emph{nāpās thavũ}) is built from \emph{nāpās}, \textquotedblleft{}not passed\textquotedblright{}, with the verb \emph{thavũ}, \textquotedblleft{}to become\textquotedblright{}. \emph{Nāpās} joins \emph{nā}, \textquotedblleft{}not\textquotedblright{}, usually traced to Persian, to a loan from English \emph{pass}.
\end{cousinweb}

\begin{grammarlens}[title={What you've built}]
One more word for everyday life at A2.
\end{grammarlens}

\subsection*{Your turn}
\begin{itemize}
\raggedright
  \item \emph{Say it:} \emph{nāpās thavũ}
  \item \emph{Say it:} \emph{nāpās thavũ}, once more
  \item \emph{Say it:} \emph{te gaṇitmā̃ nāpās thāy che} — He fails in maths.
  \item \emph{Recall:} say the Gujarati for result, then the Gujarati for marks, then say \emph{nāpās thavũ} again
\end{itemize}

\begin{tcolorbox}[breakable,colback=teal!4,colframe=teal!35!black,title={Before you move on}]
What does \gu{નાપાસ} \gu{થવું} mean? (\textquotedblleft{}To fail an exam\textquotedblright{}.) Say it once more.
\end{tcolorbox}
\section[\texorpdfstring{\gu{અભ્યાસ} (abhyās)}{અભ્યાસ (abhyās)}]{\gu{અભ્યાસ} (abhyās) — studies, study}

\label{lesson:GU-C246-abhyas}



\begin{quote}
Before the new one: say the Gujarati for marks, then the Gujarati for to fail an exam.
\end{quote}

\subsection*{You'll want to know: \gu{અભ્યાસ}}
\textbf{\gu{અભ્યાસ}} — \emph{abhyās} — \textquotedblleft{}studies, study\textquotedblright{}.

\begin{quote}
\textbf{\gu{તે} \gu{રોજ} \gu{બે} \gu{કલાક} \gu{અભ્યાસ} \gu{કરે} \gu{છે}.} — \emph{te roj be kalāk abhyās kare che} — \textquotedblleft{}She studies for two hours every day.\textquotedblright{}
\end{quote}

A masculine noun.

\begin{cousinweb}
\textbf{\gu{અભ્યાસ}} (\emph{abhyās}) is Sanskrit \emph{abhyāsa}, \textquotedblleft{}repetition, practice\textquotedblright{}, taken over whole as a tatsama. CDIAL lists it under the verb \emph{abhyasyati}, \textquotedblleft{}practises\textquotedblright{}, with Old Gujarati \emph{ahīāsivaũ}, \textquotedblleft{}to practise\textquotedblright{} (CDIAL 546).
\end{cousinweb}

\begin{grammarlens}[title={What you've built}]
One more word for everyday life at A2.
\end{grammarlens}

\subsection*{Your turn}
\begin{itemize}
\raggedright
  \item \emph{Say it:} \emph{abhyās}
  \item \emph{Say it:} \emph{abhyās}, once more
  \item \emph{Say it:} \emph{te roj be kalāk abhyās kare che} — She studies for two hours every day.
  \item \emph{Recall:} say the Gujarati for marks, then the Gujarati for to fail an exam, then say \emph{abhyās} again
\end{itemize}

\begin{tcolorbox}[breakable,colback=teal!4,colframe=teal!35!black,title={Before you move on}]
What does \gu{અભ્યાસ} mean? (\textquotedblleft{}Studies, study\textquotedblright{}.) Say it once more.
\end{tcolorbox}
\section[\texorpdfstring{\gu{મહેનત} \gu{કરવી} (mahenat karvī)}{મહેનત કરવી (mahenat karvī)}]{\gu{મહેનત} \gu{કરવી} (mahenat karvī) — to work hard}

\label{lesson:GU-C246-mahenat-karvi}



\begin{quote}
Before the new one: say the Gujarati for to fail an exam, then the Gujarati for studies.
\end{quote}

\subsection*{You'll want to know: \gu{મહેનત} \gu{કરવી}}
\textbf{\gu{મહેનત} \gu{કરવી}} — \emph{mahenat karvī} — \textquotedblleft{}to work hard\textquotedblright{}.

\begin{quote}
\textbf{\gu{અમે} \gu{પરિણામ} \gu{માટે} \gu{મહેનત} \gu{કરીએ} \gu{છીએ}.} — \emph{ame pariṇām māṭe mahenat karīe chīe} — \textquotedblleft{}We work hard for the results.\textquotedblright{}
\end{quote}

\begin{cousinweb}
\textbf{\gu{મહેનત} \gu{કરવી}} (\emph{mahenat karvī}) is built from \emph{mahenat}, \textquotedblleft{}hard work\textquotedblright{}, a loan from Arabic \emph{miḥnat}, \textquotedblleft{}trial, affliction\textquotedblright{}, which reached Gujarati through Persian (Lane), and the verb \emph{karvī}, \textquotedblleft{}to do\textquotedblright{}.
\end{cousinweb}

\begin{grammarlens}[title={What you've built}]
One more word for everyday life at A2.
\end{grammarlens}

\subsection*{Your turn}
\begin{itemize}
\raggedright
  \item \emph{Say it:} \emph{mahenat karvī}
  \item \emph{Say it:} \emph{mahenat karvī}, once more
  \item \emph{Say it:} \emph{ame pariṇām māṭe mahenat karīe chīe} — We work hard for the results.
  \item \emph{Recall:} say the Gujarati for to fail an exam, then the Gujarati for studies, then say \emph{mahenat karvī} again
\end{itemize}

\begin{tcolorbox}[breakable,colback=teal!4,colframe=teal!35!black,title={Before you move on}]
What does \gu{મહેનત} \gu{કરવી} mean? (\textquotedblleft{}To work hard\textquotedblright{}.) Say it once more.
\end{tcolorbox}
